\documentclass[11pt,a4paper]{article}
\usepackage[margin=23mm]{geometry}
\usepackage{amsmath,amssymb,amsthm}
\setlength{\emergencystretch}{2em}
\newtheorem{theorem}{Theorem}
\newcommand{\gen}[1]{\langle#1\rangle}
\newcommand{\Alt}{\operatorname{Alt}}
\title{Eight fixed points in the odd-degree orbit budget}
\author{Human Compression project: L044 continuation}
\date{}
\begin{document}\maketitle
Let $G=\gen{\rho_0,\ldots,\rho_{r-1}}\le\Alt(\Omega)$
be transitive with a connected string C-marking and all
maximal omissions intransitive. Assume $r\ge6$ and
$n=2r+1$. Retain the actual six-point head on
$D=\{1,\ldots,6\}$ and the literal interface
\[
 \tau=\rho_3=(6\ b)\delta,\quad
 R=G_{\ge3}=\Alt(\Omega\setminus\{1,\ldots,5\}),\quad
 \rho_4|_D=1,\quad G_{\ge5}|_{D\cup\{b\}}=1,
\]
where $\delta$ is supported outside $D\cup\{b\}$.
Put $H=G_{\ge5}$ and $K=G_{>5}$. Suppose every tail
colour $j\ge5$ has two full ambient-eligible edges in
one actual $H$-orbit. Normalize these assignments by the
acquired whole-orbit transport, so every assigned orbit
has intransitive $K$-action. No assignment of an eligible
pair split between distinct $H$-orbits is inferred.

\begin{theorem}
Under these hypotheses, $H$ cannot have exactly eight
fixed points.
\end{theorem}
\begin{proof}
Suppose it does. For each nontrivial $H$-orbit $O_i$,
let $F_i$ be its assigned colours and put
\[
 d_i=|F_i|-|O_i|/2.
\]
Let $p$ index the orbit assigned colour $5$. Since the
sets $F_i$ partition all $r-5$ tail colours, the equation
$n=2r+1$ gives
\[
                    \sum_i d_i=(r-5)-(n-8)/2=-3/2.       \tag{1}
\]
The acquired normalized assignment theorem gives
\[
 d_p\le0,\qquad
 d_i\le-1\text{ if }|O_i|\text{ is even},\quad
 d_i\le-\tfrac12\text{ if }|O_i|\text{ is odd}
 \quad(i\ne p).                                        \tag{2}
\]
An empty assignment has the exact deficit $-|O_i|/2$.

The odd involution $\delta$ centralizes $H$. It fixes
the seven specified fixed points and therefore also the
only eighth fixed point. We first show it preserves every
nontrivial $H$-orbit. If it exchanges orbits $O_i,O_j$
of equal size $m$, transport the assigned edges of $O_j$
by the actual involution $\tau\in G_\ell$ for every
tail colour $\ell$. On these orbits its action is that
of $\delta$, and it centralizes $H$. Thus each colour in
$F_i\cup F_j$ has two ambient-eligible edges in $O_i$.
They are also eligible for the full local omission of
$H|_{O_i}$, by subgroup inclusion. The complete selected
colour bound gives
\[
 |F_i|+|F_j|\le\lfloor m/2\rfloor,
 \qquad d_i+d_j\le-\lceil m/2\rceil.
\]
For $m\ge3$ this is at most $-2$, contrary to (1)--(2).
For $m=2$, neither orbit can contain two disjoint edges
of one involution, so both assigned sets are empty and
the deficit is exactly $-2$, again impossible.

Because $\delta$ is odd and fixes the eight fixed points,
it acts oddly on some preserved nontrivial orbit $O$.
Its restriction there is a nonidentity involution
centralizing a transitive group, hence is free.
Write $|O|=2q$, where $q$ is odd.

First suppose $q\ge3$. Its assignment cannot be empty,
which would contribute $-q\le-3$ to (1). The acquired
internal-pair forest gives $d_O\le-1$. If $O\ne O_p$,
integrality and (1)--(2) force $d_O=-1$.
The reviewed odd-degree row theorem applies: its assigned
$q-1$ colours form a consecutive interval whose last
label $z$ is strictly below $r-1$, and
\[
                  |\operatorname{Fix}(G_{>z})|\ge2z+3.
\]
Every colour after $z$ retains its ambient eligible pair.
The full terminal interval bound therefore gives
\[
 |\operatorname{supp}(G_{>z})|
                  \ge2(r-1-z)+1.
\]
The interval is nonempty and is intransitive on $\Omega$
because it fixes the specified head. Adding the last
two inequalities yields $n\ge2r+2$, contrary to $n=2r+1$.

If instead $O=O_p$, (1)--(2) and integrality give
$d_p=-1$. The remaining total deficit is $-1/2$.
There is therefore exactly one further nontrivial orbit
$V$, of odd size, with $d_V=-1/2$. In particular $F_V$
is nonempty; a nontrivial odd orbit with empty assignment
would contribute at most $-3/2$. The internal equality
theorem gives
\[
 F_p=\{5,\ldots,q+3\},\qquad
 I=F_V=\{q+4,\ldots,r-1\},\quad 2|I|=|V|-1,
\]
and every colour outside $F_p$ has trivial image on the
$\delta$-pairs of $O$. Every possible edge of a colour
in $I$ on $O$ is consequently one of these pairs, already
joined by $\tau\in G_j$, and is ambient ineligible.
The eight fixed points carry no tail edges. All eligible
edges of colours in $I$ are therefore confined to $V$.
Since $T=G_I\le K$ and $K$ is intransitive on $V$, the
acquired localization theorem gives
\[
 |V|-1=2|I|\le |V|-o_T(V)\le |V|-2,
\]
a contradiction. This excludes $q\ge3$ in every position.

Finally $q=1$. No tail edge on the two-point orbit $O$
is eligible, since its points are joined by $\tau$ in
every tail omission. Its assignment is empty and
$d_O=-1$. Only $-1/2$ of the total deficit remains.
The involution $\delta$ cannot act nontrivially on any
other nontrivial orbit $W$. Such an action would be free.
For $|W|=2$, the same argument contributes another $-1$.
For $|W|\ge4$, an empty assignment contributes at most
$-2$; a nonempty assignment contributes at most $-1$
by the internal pair forest. Its elementary forest bound
also holds for two pairs: at most one distinct selected
colour can join the two pair vertices. Each possibility
contradicts (1)--(2), including when $W=O_p$.

Thus $\delta$ is literally the one transposition on $O$.
The reviewed odd-degree exterior-pair obstruction applies:
all its actual head and natural-$R$ hypotheses have been
retained, and every colour $j\ge6$ has an ambient eligible
pair by the assignment. It excludes precisely this
$\tau=(6\ b)\delta$ configuration for $n=2r+1$.
The final case is impossible, proving the theorem.
\end{proof}

\paragraph{Attribution and remaining scope.}
The normalized assignment and selected-colour inequalities,
internal pair equality, complete localization proof, strict
odd-degree row exit and exact exterior-pair obstruction are
acquired inputs. The new step assembles their full parity and
deficit cases at exactly eight fixed points, including exchanged
orbits and an empty internal pair. This does not supply the
assignment when eligible pairs split across different orbits.
The seven-fixed-point budgets and the global separator argument
remain. No alternating rank induction, new computation, formal
declaration, original or question credit is claimed.
\end{document}
